Table~\ref{tab:attention} compares what the core model and the limited-attention model predict for the non-disclosure premium, and reports the attention estimates. We estimate $1-\psi_\delta$ as the ratio of the cross-market sums of the two sides of~\eqref{eq:att_lambda}. Absent selection, and unlike a regression through the origin, this estimator is free of the attenuation and correlated-noise biases that arise when $\hat\beta_{1,m}$ is noisy and shares the E-to-G index with $\hat\Delta_m$. The estimation sample is, however, selected on $\hat\beta_{1,m}>0$, and selecting on a noisy $\hat\beta_{1,m}$ biases both estimators toward full attention: a Monte Carlo calibrated to our markets puts the ratio of sums between \attMCromSelLo{} and \attMCromSelHi{} for a true value of \attMCtrue{} (Appendix~\ref{app:att_estimator}). The apartment estimate is therefore, if anything, an upper bound on attention to silence, and the house estimate may overstate it. We estimate one $\psi_\delta$ per segment rather than per market. A market-specific value is identified by the market's own ratio $\hat\Delta_m/\hat\beta_{1,m}$, but it would fit each premium exactly, place every market in Partial Pooling by construction, and leave no restriction to test, and the market ratios are too noisy to be informative on their own.\footnote{The estimators rank by how much weight they place on markets with a large energy-efficiency premium: the median of market ratios gives $\attPsiMedApt$, the ratio of sums $\attPsiApt$, the regression through the origin $\attPsiOLSApt$, and a regression with a constant, which the model excludes, $\attPsiConstApt$ (intercept $\attIcptApt$, s.e.\ $\attIcptSeApt$).}
